\chapter{Einleitung}

\section{Versuchsziel}

Eine fundamentale Erkenntnis der modernen Physik ist die Quantisierung der möglichen Energieniveaus von Atomen und ihren Kernen und Hüllenelektronen \cite[112]{Demtroeder3}. Mit dem \textit{Franck-Hertz}-Versuchs wird die Energiequantisierung von Atomenergiezuständen bei Quecksilber-Atomen nachgewiesen werden. Außerdem wird anhand der Beobachtung der
Energieaufspaltung durch den normalen Zeeman-Effekts für die rote Cadmium (Cd)-Spektrallinie mit Wellenlänge $\lambda_0=\SI{643.847}{\nano\meter}$ \cite{NISTcd} bei einem externen, homogenen Magnetfeld der Wert des Bohr'schen Magnetons $\mu_B$ abhängig von der Magnetfeldstärke bestimmt.

\subsection{Franck-Hertz-Versuch}

Der experimentelle Nachweis der diskreten Energieniveaus von Atomen wurde 1911-1914 von \uppercase{Franck} und \uppercase{Hertz} mit dem \textit{Franck-Hertz-Versuch} erbracht. Hierbei werden Elektronen durch eine angelegte Spannung beschleunigt und durchlaufen dann eine Röhre, welche mit einem Gas unter niedrigem Druck -- in der durchgeführten Variante wird hierfür Quecksilber (Hg) genutzt -- gefüllt ist, von einer Kathode zu einer Anode. Erreichen die Elektronen die andere Seite der Röhre, werden sie als Anodenstrom gemessen. Auf dem Weg stoßen die Elektronen solange nur inelastisch mit Quecksilberatomen bis ihre durch die Beschleunigungsspannung erreichte Energie der Anregungsenergie zwischen dem nächsthöheren Hg-Energiezustand entspricht. Diese Energie ist charakteristisch für das verwendete Gas, den Ausgangszustand der Gasatome sowie den angeregten Übergang, d.h. auch den finalen Atomzustand. Es ist zudem noch ein Gegenfeld angelegt, sodass sehr langsame Elektronen, die nach einem Stoß fast keine kinetische Energie mehr besitzen, nicht die Anode erreichen können.
Wenn die Elektronen gerade genug Energie besitzen um ein Atom in den höheren Zustand anzuregen, haben sie danach nicht mehr genug kinetische Energie um trotz dem Gegenfeld die Anode zu erreichen und der gemessene Strom sinkt für diese Beschleunigungsspannung (d.h. für diese kinetische Energie der Elektronen) stark ab. Es zeigt sich im Anodenstrom also ein Resonanzmuster für die Anregungsenergie der Gasatome (für den Hüllenelektronen-Übergang, dessen Energiedifferenz mit den genutzten angelegten Spannungen). Ist die Elektronenenergie geringer als die Anregungsenergie, können keine elastischen Stöße stattfinden und der Anodenstrom zeigt keinen Abfall. Wird durch die Beschleunigungsspannung ein Vielfaches der Anregungsenergie erreicht, können die Elektronen auf dem Weg zur Anode zweimal mit Gasatomen elastisch stoßen, sodass auch bei Vielfachen der Anregungsenergie in der Spannung ein Resonanzabfall im Anodenstom auftritt. \cite[687-688]{gerthsen}.

\todo{Aufbau und so hier rein}

\subsection{Normaler Zeeman-Effekt}


Die Quantisierung der Atomzustände werden durch Quantenzahlen beschrieben. Unterscheiden sich die Quantenzahlen zweier Zustände, unterscheiden sich auch die Energien. Die Energiespektren der Zustände eines Ein- oder Mehrelektronenatoms können über Termschemata dargestellt werden \cite[1]{Versuchsanleitung401}.
